\documentclass[11pt,a4paper]{article}
\usepackage{amsmath}
\usepackage{amssymb}
\usepackage[table]{xcolor}
\usepackage{tabularx}
\usepackage{multirow}
\usepackage{geometry}
\usepackage{enumitem}
\usepackage{parskip}
\usepackage{titlesec}
\usepackage{tocloft}
\usepackage[most]{tcolorbox}

\geometry{margin=2.5cm}

\titleformat{\section}{\large\bfseries}{}{0em}{}
\setlength{\cftbeforesecskip}{6pt}

% ===== Pastel colors, one per question (cycled across 40 questions) =====
\definecolor{q1color}{HTML}{D6EAF8} % pastel blue
\definecolor{q2color}{HTML}{D5F5E3} % pastel green
\definecolor{q3color}{HTML}{FDEBD0} % pastel orange
\definecolor{q4color}{HTML}{F9E79F} % pastel yellow
\definecolor{q5color}{HTML}{E8DAEF} % pastel purple
\definecolor{q6color}{HTML}{FADBD8} % pastel red/pink
\definecolor{q7color}{HTML}{D1F2EB} % pastel teal
\definecolor{q8color}{HTML}{FCF3CF} % pastel cream
\definecolor{q9color}{HTML}{EBDEF0} % pastel lavender

% ===== Mark-scheme table (PDF-style grid, tinted per question) =====
% Retained for preamble compatibility with the Paper 2 file so the two can be
% merged without editing. Not used in this multiple-choice paper.
\newenvironment{mstab}[1]{%
  {\color{#1!75!black}\nobreak Mark Scheme}\par\nopagebreak\smallskip
  \arrayrulecolor{#1!70!black}%
  \renewcommand{\arraystretch}{1.3}%
  \tabularx{\textwidth}{%
    |>{\columncolor{#1!12}\raggedright\arraybackslash}p{1.9cm}|%
     >{\columncolor{#1!12}}X|%
     >{\columncolor{#1!12}\centering\arraybackslash}p{1.1cm}|}%
  \hline
  \rowcolor{#1!45} Question & Answer & Marks \\ \hline
}{%
  \endtabularx
  \arrayrulecolor{black}%
  \par\medskip
}

% ===== Answer colorbox style =====
\newtcolorbox{ansbox}[1]{
  colback=#1!8,
  colbacktitle=#1!30,
  colframe=#1!70!black,
  coltitle=black,
  fonttitle=\normalfont,
  title=Worked Solution,
  boxrule=0.6pt,
  arc=2mm,
  left=4pt, right=4pt, top=4pt, bottom=4pt,
  breakable
}

% ===== Topic label for question headings (right-aligned, subtle) =====
\newcommand{\topiclabel}[1]{\hfill\normalfont\small\itshape\textcolor{black!55}{#1}}

% Per-question head: \examq{paper-id}{number}{topics}{marks}
% Same command and rendered line as the question paper.
\newcommand{\examq}[4]{\section[#1 - Q#2]{#1 - Q#2 - #3 - #4}}

\begin{document}

\begin{center}
  {\LARGE \textbf{Cambridge O Level Physics 5054/11}}\\[6pt]
  {\large May/June 2025 -- Multiple Choice: Answers \& Solutions}
\end{center}

\bigskip

%%%%%%%%%%%%%%%%%%%
\examq{5054/11/M/J/25}{1}{Physical Quantities and Measurements}{1}
%%%%%%%%%%%%%%%%%%%

\begin{ansbox}{q1color}
Answer: $\boxed{\text{D}}$\\[4pt]
Weight is a force (the pull of gravity), so it has both a magnitude and a direction, making it a vector. Mass has magnitude only and is a scalar, so the correct pairing is weight, vector.
\end{ansbox}

%%%%%%%%%%%%%%%%%%%
\examq{5054/11/M/J/25}{2}{Motion or Kinematics}{1}
%%%%%%%%%%%%%%%%%%%

\begin{ansbox}{q2color}
Answer: $\boxed{\text{B}}$\\[4pt]
The distance travelled is the area under the speed--time graph. Splitting the region beneath the line into triangular and trapezoidal sections and adding their areas over the $10\,\mathrm{s}$ gives
\[
s = \boxed{25\ \mathrm{m}}
\]
\end{ansbox}

%%%%%%%%%%%%%%%%%%%
\examq{5054/11/M/J/25}{3}{Forces or Dynamics}{1}
%%%%%%%%%%%%%%%%%%%

\begin{ansbox}{q3color}
Answer: $\boxed{\text{D}}$\\[4pt]
Terminal velocity is only reached once the air resistance has grown large enough to equal the weight. Over a short fall the coin has not sped up enough for this to happen, so its weight is always greater than the air resistance -- there is still a resultant downward force, so it keeps accelerating and never reaches terminal velocity.
\end{ansbox}

%%%%%%%%%%%%%%%%%%%
\examq{5054/11/M/J/25}{4}{Deformation}{1}
%%%%%%%%%%%%%%%%%%%

\begin{ansbox}{q4color}
Answer: $\boxed{\text{B}}$\\[4pt]
Point X is where the load--extension line stops being straight, i.e.\ where extension is no longer proportional to load. This is called the limit of proportionality.
\end{ansbox}

%%%%%%%%%%%%%%%%%%%
\examq{5054/11/M/J/25}{5}{Mass Weight and Density}{1}
%%%%%%%%%%%%%%%%%%%

\begin{ansbox}{q5color}
Answer: $\boxed{\text{B}}$\\[4pt]
All the samples are the same type of wood, so they share one density. Since $\rho = m/V$ is constant, mass is directly proportional to volume, giving a straight line that passes through the origin -- graph B.
\end{ansbox}

%%%%%%%%%%%%%%%%%%%
\examq{5054/11/M/J/25}{6}{Forces or Dynamics}{1}
%%%%%%%%%%%%%%%%%%%

\begin{ansbox}{q6color}
Answer: $\boxed{\text{D}}$\\[4pt]
An object moving in a circle at constant speed still has a resultant force, and that force (the centripetal force) always points towards the centre of the circle. Arrow D is the one directed from the truck towards the centre of the vertical loop.
\end{ansbox}

%%%%%%%%%%%%%%%%%%%
\examq{5054/11/M/J/25}{7}{Momentum (Dynamics)}{1}
%%%%%%%%%%%%%%%%%%%

\begin{ansbox}{q7color}
Answer: $\boxed{\text{A}}$\\[4pt]
Momentum is $p = mv$, so its unit is (mass)$\times$(velocity) $=$ kilogram $\times$ metre per second:
\[
\boxed{\mathrm{kg\,m/s}}
\]
\end{ansbox}

%%%%%%%%%%%%%%%%%%%
\examq{5054/11/M/J/25}{8}{Momentum (Dynamics)}{1}
%%%%%%%%%%%%%%%%%%%

\begin{ansbox}{q8color}
Answer: $\boxed{\text{C}}$\\[4pt]
Impulse equals force $\times$ time, so the force is impulse $\div$ time. Converting $70\,\mathrm{ms} = 0.070\,\mathrm{s}$:
\[
F = \frac{\text{impulse}}{t} = \frac{14\,000}{0.070} = \boxed{200\,000\ \mathrm{N}}
\]
\end{ansbox}

%%%%%%%%%%%%%%%%%%%
\examq{5054/11/M/J/25}{9}{Work Energy and Power}{1}
%%%%%%%%%%%%%%%%%%%

\begin{ansbox}{q9color}
Answer: $\boxed{\text{C}}$\\[4pt]
Work done $=$ force $\times$ distance moved in the direction of that force. The weight and the contact force act vertically, perpendicular to the horizontal motion, so they do no work. Friction acts backwards and removes kinetic energy. It is the forward pulling force whose work increases the kinetic energy store.
\end{ansbox}

%%%%%%%%%%%%%%%%%%%
\examq{5054/11/M/J/25}{10}{Mass Weight and Density}{1}
%%%%%%%%%%%%%%%%%%%

\begin{ansbox}{q1color}
Answer: $\boxed{\text{A}}$\\[4pt]
Mass is a measure of the amount of matter in an object and also of its inertia (its resistance to a change in motion). Since P has the greater mass, P has both the greater amount of matter and the greater resistance to change of motion.
\end{ansbox}

%%%%%%%%%%%%%%%%%%%
\examq{5054/11/M/J/25}{11}{Pressure}{1}
%%%%%%%%%%%%%%%%%%%

\begin{ansbox}{q2color}
Answer: $\boxed{\text{C}}$\\[4pt]
The tiles are identical, so each has the same weight and therefore exerts the same force on the earth. The vertical tile rests on a much smaller contact area, and since pressure $=$ force $\div$ area, it exerts a greater pressure -- which is why it sinks. So the forces are the same but the pressures are different.
\end{ansbox}

%%%%%%%%%%%%%%%%%%%
\examq{5054/11/M/J/25}{12}{Pressure}{1}
%%%%%%%%%%%%%%%%%%%

\begin{ansbox}{q3color}
Answer: $\boxed{\text{C}}$\\[4pt]
Pressure in a liquid increases with depth below the surface ($p = \rho g h$) and does not depend on the shape of the container. Point C is the deepest below the liquid surface, so the pressure there is greatest.
\end{ansbox}

%%%%%%%%%%%%%%%%%%%
\examq{5054/11/M/J/25}{13}{Thermal Properties \& Temperature}{1}
%%%%%%%%%%%%%%%%%%%

\begin{ansbox}{q4color}
Answer: $\boxed{\text{A}}$\\[4pt]
The particles in a solid are held by strong forces, so heating does not push them very far apart. This is why solids expand less than liquids and gases (whose particle forces are weaker).
\end{ansbox}

%%%%%%%%%%%%%%%%%%%
\examq{5054/11/M/J/25}{14}{Kinetic Model of Matter}{1}
%%%%%%%%%%%%%%%%%%%

\begin{ansbox}{q5color}
Answer: $\boxed{\text{C}}$\\[4pt]
Heating gives the air molecules more kinetic energy, so they move faster. In the fixed volume they therefore strike the walls more often (and harder), and more frequent collisions with the walls mean a greater pressure. The molecules themselves do not expand, and their number stays the same.
\end{ansbox}

%%%%%%%%%%%%%%%%%%%
\examq{5054/11/M/J/25}{15}{Thermal Properties \& Temperature}{1}
%%%%%%%%%%%%%%%%%%%

\begin{ansbox}{q6color}
Answer: $\boxed{\text{A}}$\\[4pt]
To find $c$ from $E = mc\,\Delta\theta$ the student needs the mass (balance, 1), the temperature rise (thermometer, 3) and the energy supplied. As the heater's power is already known ($50\,\mathrm{W}$), the energy is power $\times$ time, so only the heater (4) and a stopwatch (5) are needed -- there is no need for the ammeter and voltmeter. The apparatus is 1, 3, 4 and 5.
\end{ansbox}

%%%%%%%%%%%%%%%%%%%
\examq{5054/11/M/J/25}{16}{Thermal Properties \& Temperature}{1}
%%%%%%%%%%%%%%%%%%%

\begin{ansbox}{q7color}
Answer: $\boxed{\text{C}}$\\[4pt]
Only the fastest, most energetic molecules have enough energy to escape from the surface. Losing these high-energy molecules lowers the average kinetic energy of those left behind, so the liquid cools.
\end{ansbox}

%%%%%%%%%%%%%%%%%%%
\examq{5054/11/M/J/25}{17}{Thermal Properties \& Temperature}{1}
%%%%%%%%%%%%%%%%%%%

\begin{ansbox}{q8color}
Answer: $\boxed{\text{C}}$\\[4pt]
Evaporation is fastest with a high temperature, a large surface area and a high air speed across the surface. Row C combines the highest temperature ($38\,^{\circ}$C), the largest area ($70\,\mathrm{cm^2}$) and the greatest air speed ($5.0\,\mathrm{m/s}$), so it evaporates most rapidly.
\end{ansbox}

%%%%%%%%%%%%%%%%%%%
\examq{5054/11/M/J/25}{18}{Transfer of Thermal Energy}{1}
%%%%%%%%%%%%%%%%%%%

\begin{ansbox}{q9color}
Answer: $\boxed{\text{B}}$\\[4pt]
Metals conduct thermal energy so well because they contain free (delocalised) electrons. These gain energy in the hot region and move quickly through the metal, carrying that energy to the ions in the cooler region. (The electrons are free to move, not fixed, and it is electrons -- not protons or ions -- that travel.)
\end{ansbox}

%%%%%%%%%%%%%%%%%%%
\examq{5054/11/M/J/25}{19}{Light (Reflection/Refraction/Lenses)}{1}
%%%%%%%%%%%%%%%%%%%

\begin{ansbox}{q1color}
Answer: $\boxed{\text{B}}$\\[4pt]
The angle of incidence is always measured between the incident (incoming) ray and the normal, not the mirror surface. That is angle B.
\end{ansbox}

%%%%%%%%%%%%%%%%%%%
\examq{5054/11/M/J/25}{20}{Light (Reflection/Refraction/Lenses)}{1}
%%%%%%%%%%%%%%%%%%%

\begin{ansbox}{q2color}
Answer: $\boxed{\text{B}}$\\[4pt]
The ray from the shoes to the eyes reflects off the mirror at a point exactly halfway between the eye level ($150\,\mathrm{cm}$) and the shoes ($0\,\mathrm{cm}$), i.e.\ at $75\,\mathrm{cm}$. This reflecting point must lie on the mirror. With a $50\,\mathrm{cm}$ mirror whose top is at $h + 50$, we need
\[
h + 50 \geq 75 \quad\Rightarrow\quad h \geq 25\ \mathrm{cm}
\]
so the smallest height is $\boxed{25\ \mathrm{cm}}$.
\end{ansbox}

%%%%%%%%%%%%%%%%%%%
\examq{5054/11/M/J/25}{21}{Transfer of Thermal Energy}{1}
%%%%%%%%%%%%%%%%%%%

\begin{ansbox}{q3color}
Answer: $\boxed{\text{C}}$\\[4pt]
A dull black surface is the best possible surface for radiation: it is a good emitter and a good absorber of infrared, and (being dark and dull) a poor reflector. So: good emitter, good absorber and poor reflector.
\end{ansbox}

%%%%%%%%%%%%%%%%%%%
\examq{5054/11/M/J/25}{22}{Light (Reflection/Refraction/Lenses)}{1}
%%%%%%%%%%%%%%%%%%%

\begin{ansbox}{q4color}
Answer: $\boxed{\text{B}}$\\[4pt]
The type of each lens is read from how it bends the rays. A converging lens brings the rays together (inwards); a diverging lens spreads them apart (outwards). At positions 1 and 2 the rays are brought inwards (converging lenses), while at position 3 the rays are spread outwards (a diverging lens): converging, converging, diverging.
\end{ansbox}

%%%%%%%%%%%%%%%%%%%
\examq{5054/11/M/J/25}{23}{Light (Reflection/Refraction/Lenses)}{1}
%%%%%%%%%%%%%%%%%%%

\begin{ansbox}{q5color}
Answer: $\boxed{\text{C}}$\\[4pt]
The ray passes from the optically denser water into the less dense air, so at the surface it refracts away from the normal (bending towards the water surface). Path C is the one that shows the ray bending away from the normal as it leaves the water.
\end{ansbox}

%%%%%%%%%%%%%%%%%%%
\examq{5054/11/M/J/25}{24}{Light (Reflection/Refraction/Lenses)}{1}
%%%%%%%%%%%%%%%%%%%

\begin{ansbox}{q6color}
Answer: $\boxed{\text{D}}$\\[4pt]
As a magnifying glass, the object is placed inside the focal length of the converging lens. The image formed is then virtual, upright and magnified -- so its nature is virtual and upright.
\end{ansbox}

%%%%%%%%%%%%%%%%%%%
\examq{5054/11/M/J/25}{25}{Sound}{1}
%%%%%%%%%%%%%%%%%%%

\begin{ansbox}{q7color}
Answer: $\boxed{\text{A}}$\\[4pt]
Sound is a longitudinal wave (the particles vibrate along the direction the wave travels). A rarefaction is where the particles are spread out, so the pressure there is smaller than in a compression: longitudinal, smaller.
\end{ansbox}

%%%%%%%%%%%%%%%%%%%
\examq{5054/11/M/J/25}{26}{Sound}{1}
%%%%%%%%%%%%%%%%%%%

\begin{ansbox}{q8color}
Answer: $\boxed{\text{A}}$\\[4pt]
Loudness depends on the amplitude of the sound wave. A louder note means a larger amplitude. (Frequency changes pitch, not loudness, and speed and wavelength do not set how loud the note is.)
\end{ansbox}

%%%%%%%%%%%%%%%%%%%
\examq{5054/11/M/J/25}{27}{Sound}{1}
%%%%%%%%%%%%%%%%%%%

\begin{ansbox}{q9color}
Answer: $\boxed{\text{A}}$\\[4pt]
The extra time for the pulse that travels on to the bottom surface is
\[
\Delta t = 1.5 \times 10^{-5} - 1.0 \times 10^{-5} = 0.5 \times 10^{-5}\ \mathrm{s}
\]
This is the time for the pulse to travel the extra distance down to the bottom \emph{and back}, i.e.\ $2x$:
\[
x = \frac{v\,\Delta t}{2} = \frac{5200 \times 0.5 \times 10^{-5}}{2} = \boxed{0.013\ \mathrm{m}}
\]
\end{ansbox}

%%%%%%%%%%%%%%%%%%%
\examq{5054/11/M/J/25}{28}{Electrical Quantities}{1}
%%%%%%%%%%%%%%%%%%%

\begin{ansbox}{q1color}
Answer: $\boxed{\text{B}}$\\[4pt]
The current is $I = V/R = 12/100 = 0.12\,\mathrm{A}$, and $30\,\text{min} = 1800\,\mathrm{s}$. The charge is
\[
Q = It = 0.12 \times 1800 = 216\ \mathrm{C} \approx \boxed{220\ \mathrm{C}}
\]
\end{ansbox}

%%%%%%%%%%%%%%%%%%%
\examq{5054/11/M/J/25}{29}{Electrical Quantities}{1}
%%%%%%%%%%%%%%%%%%%

\begin{ansbox}{q2color}
Answer: $\boxed{\text{C}}$\\[4pt]
Conventional current flows around the external circuit from the positive terminal to the negative terminal. The free electrons, being negatively charged, drift the opposite way -- from negative to positive. So the row is: conventional current positive$\rightarrow$negative, electrons negative$\rightarrow$positive.
\end{ansbox}

%%%%%%%%%%%%%%%%%%%
\examq{5054/11/M/J/25}{30}{Electric Circuits}{1}
%%%%%%%%%%%%%%%%%%%

\begin{ansbox}{q3color}
Answer: $\boxed{\text{D}}$\\[4pt]
Reading the graph, at $10\,^{\circ}$C the thermistor's resistance is $80\,\Omega$. In series the resistances simply add:
\[
R_\text{total} = 80 + 20 = \boxed{100\ \Omega}
\]
\end{ansbox}

%%%%%%%%%%%%%%%%%%%
\examq{5054/11/M/J/25}{31}{Electric Circuits}{1}
%%%%%%%%%%%%%%%%%%%

\begin{ansbox}{q4color}
Answer: $\boxed{\text{C}}$\\[4pt]
Two of the $6.0\,\Omega$ resistors are in parallel:
\[
\frac{1}{R_p} = \frac{1}{6.0} + \frac{1}{6.0} \quad\Rightarrow\quad R_p = 3.0\ \Omega
\]
This combination is in series with the third resistor:
\[
R_\text{total} = 3.0 + 6.0 = \boxed{9.0\ \Omega}
\]
\end{ansbox}

%%%%%%%%%%%%%%%%%%%
\examq{5054/11/M/J/25}{32}{Electro Magnetism / Electromagnetic Induction}{1}
%%%%%%%%%%%%%%%%%%%

\begin{ansbox}{q5color}
Answer: $\boxed{\text{C}}$\\[4pt]
The induced e.m.f.\ is greatest when the plane of the coil is parallel to the field ($0^{\circ}$ or $180^{\circ}$), because the coil sides then cut field lines fastest, and it is zero when the plane is perpendicular to the field ($90^{\circ}$). Matching the trace, point 1 sits at a zero ($90^{\circ}$), point 2 at a peak ($180^{\circ}$) and point 3 at a peak ($0^{\circ}$).
\end{ansbox}

%%%%%%%%%%%%%%%%%%%
\examq{5054/11/M/J/25}{33}{Electro Magnetism / Electromagnetic Induction}{1}
%%%%%%%%%%%%%%%%%%%

\begin{ansbox}{q6color}
Answer: $\boxed{\text{C}}$\\[4pt]
Beta particles are negative, so a downward beam of them is equivalent to a conventional current flowing upwards. With the field pointing from the N pole across to the S pole and the (conventional) current upwards, Fleming's left-hand rule gives a force directed into the page.
\end{ansbox}

%%%%%%%%%%%%%%%%%%%
\examq{5054/11/M/J/25}{34}{Electro Magnetism / Electromagnetic Induction}{1}
%%%%%%%%%%%%%%%%%%%

\begin{ansbox}{q7color}
Answer: $\boxed{\text{B}}$\\[4pt]
A step-up transformer raises the voltage, which lowers the current in the transmission cable (since $P = VI$). A smaller current means far less energy wasted as heat in the cable. Transformers work only with a changing magnetic field, so they operate on a.c. So: to decrease the current in the cable, using a.c.
\end{ansbox}

%%%%%%%%%%%%%%%%%%%
\examq{5054/11/M/J/25}{35}{The Nuclear Atom}{1}
%%%%%%%%%%%%%%%%%%%

\begin{ansbox}{q8color}
Answer: $\boxed{\text{B}}$\\[4pt]
Isotopes of the same element have the same number of protons but different numbers of neutrons. Diagrams 1 and 4 are the pair with the same proton number but different neutron numbers.
\end{ansbox}

%%%%%%%%%%%%%%%%%%%
\examq{5054/11/M/J/25}{36}{Radioactivity}{1}
%%%%%%%%%%%%%%%%%%%

\begin{ansbox}{q9color}
Answer: $\boxed{\text{D}}$\\[4pt]
An alpha decay lowers the nucleon number by 4 and the proton number by 2; a beta decay raises the proton number by 1 and leaves the nucleon number unchanged. To end with the nucleon number down by 4 but the proton number unchanged, one alpha ($-4$ nucleons, $-2$ protons) must be balanced by two betas ($+2$ protons, $0$ nucleons): one alpha-particle and two beta-particles.
\end{ansbox}

%%%%%%%%%%%%%%%%%%%
\examq{5054/11/M/J/25}{37}{Radioactivity}{1}
%%%%%%%%%%%%%%%%%%%

\begin{ansbox}{q1color}
Answer: $\boxed{\text{A}}$\\[4pt]
The half-life is the time for the count rate to fall to half its value. Reading the graph, the count rate drops to half its starting value after $60$ minutes, so the half-life is $60$ minutes.
\end{ansbox}

%%%%%%%%%%%%%%%%%%%
\examq{5054/11/M/J/25}{38}{The Nuclear Atom}{1}
%%%%%%%%%%%%%%%%%%%

\begin{ansbox}{q2color}
Answer: $\boxed{\text{B}}$\\[4pt]
Nucleon number is conserved. Before: $235 + 1 = 236$ nucleons. After: the two daughters plus the three released neutrons. So
\[
236 = 92 + X + 3 \quad\Rightarrow\quad X = 236 - 95 = \boxed{141}
\]
\end{ansbox}

%%%%%%%%%%%%%%%%%%%
\examq{5054/11/M/J/25}{39}{Earth \& Solar System}{1}
%%%%%%%%%%%%%%%%%%%

\begin{ansbox}{q3color}
Answer: $\boxed{\text{B}}$\\[4pt]
Planets closer to the Sun orbit faster and complete an orbit in less time. Earth is closer than Jupiter, so Earth has the greater orbital speed $v$, while Jupiter, being further out, has the greater orbital period $T$.
\end{ansbox}

%%%%%%%%%%%%%%%%%%%
\examq{5054/11/M/J/25}{40}{Earth \& Solar System}{1}
%%%%%%%%%%%%%%%%%%%

\begin{ansbox}{q4color}
Answer: $\boxed{\text{C}}$\\[4pt]
The order of planets outward from the Sun begins Mercury, Venus, Earth, so Earth is the third closest planet to the Sun.
\end{ansbox}

\end{document}